\section{Recherche par Voisinage Variable (VNS)}

La méthode VNS (\textit{Variable Neighborhood Search}) est une métaheuristique qui explore plusieurs structures de voisinage pour échapper aux optima locaux. L'idée clé : un optimum local pour un voisinage ne l'est pas nécessairement pour un autre.

\subsection{Solution initiale}

On part de la solution obtenue par l'heuristique LTR, avec $\sum T_j = 22$ ($C_1 = 19$, $T_1 = 13$ ; $C_2 = 14$, $T_2 = 6$ ; $C_3 = 8$, $T_3 = 3$).

\subsection{Structures de voisinage}

On définit trois voisinages de taille croissante :

\begin{enumerate}
    \item[\textbf{$N_1$}] \textbf{Swap intra-machine :} échanger deux jobs consécutifs sur une même machine.
    \item[\textbf{$N_2$}] \textbf{Swap inter-machine :} pour un même job, modifier sa priorité sur deux machines différentes (avancer sur l'une, reculer sur l'autre).
    \item[\textbf{$N_3$}] \textbf{Réinsertion :} retirer un job d'une position sur une machine et le réinsérer à une autre position.
\end{enumerate}

\textit{Remarque :} $N_1 \subset N_3$ (un swap consécutif est un cas particulier de réinsertion). $N_2$ agit sur deux machines à la fois, contrairement à $N_1$ et $N_3$ qui sont intra-machine.

\subsection{Deux variantes}

Nous implémentons et comparons deux variantes du VNS :

\subsubsection{VNS déterministe (\textit{best improvement})}

\begin{enumerate}
    \item Poser $k = 1$.
    \item Tester \textbf{tous} les voisins dans $N_k$, garder le meilleur.
    \item Si amélioration : accepter et $k \leftarrow 1$. Sinon : $k \leftarrow k+1$.
    \item Arrêt quand $k > 3$.
\end{enumerate}

Cette variante est exhaustive : elle garantit de trouver le meilleur voisin dans chaque voisinage. Elle est déterministe (même résultat à chaque exécution).

\subsubsection{VNS stochastique (classique)}

\begin{enumerate}
    \item Poser $k = 1$.
    \item \textbf{Perturbation :} générer un voisin aléatoire $S'$ dans $N_k$.
    \item \textbf{Recherche locale :} améliorer $S'$ par descente dans $N_1$ (tant qu'un swap consécutif améliore, on l'applique) pour obtenir $S''$.
    \item Si $\sum T_j(S'') < \sum T_j(S)$ : $S \leftarrow S''$, $k \leftarrow 1$. Sinon : $k \leftarrow k+1$.
    \item Arrêt quand $k > 3$ ou nombre maximal d'itérations atteint.
\end{enumerate}

Cette variante explore plus largement l'espace de solutions grâce à l'aléatoire, ce qui est avantageux sur les grandes instances.

\subsection{Déroulement et solution}

\textit{Légende des couleurs : \textcolor{blue}{Job 1}, \textcolor{red}{Job 2}, \textcolor{green!70!black}{Job 3}.}

Partant de LTR ($\sum T_j = 22$), la variante \textbf{déterministe} améliore deux fois de suite dans $N_1$ (passage à 11 puis à 10), puis $N_1$, $N_2$ et $N_3$ ne trouvent plus d'amélioration : arrêt à $\sum T_j = 10$. La variante \textbf{stochastique} atteint $\sum T_j = 10$ dès la première perturbation suivie de recherche locale, puis stagne. Les deux variantes convergent donc vers la même solution :

\begin{figure}[H]
\centering
\begin{ganttchart}[
    vgrid, hgrid,
    x unit=1.0cm, y unit chart=0.7cm,
    bar/.append style={draw=black, line width=1pt}
]{1}{11}
\gantttitle{Temps}{11} \\
\gantttitlelist{1,...,11}{1} \\
\ganttbar[bar/.style={fill=green!40}]{$M_1$}{1}{2}
\ganttbar[bar/.style={fill=red!40}]{}{3}{3}
\ganttbar[bar/.style={fill=blue!40}]{}{4}{6} \\
\ganttbar[bar/.style={fill=green!40}]{$M_2$}{3}{3}
\ganttbar[bar/.style={fill=red!40}]{}{4}{7}
\ganttbar[bar/.style={fill=blue!40}]{}{8}{9} \\
\ganttbar[bar/.style={fill=green!40}]{$M_3$}{4}{8}
\ganttbar[bar/.style={fill=red!40}]{}{9}{10}
\ganttbar[bar/.style={fill=blue!40}]{}{11}{11}
\end{ganttchart}
\caption{Solution VNS : $C_1=11$ ($T_1=5$), $C_2=10$ ($T_2=2$), $C_3=8$ ($T_3=3$) ; $\sum T_j = 10$}
\end{figure}

\subsection{Comparaison des deux variantes}

\begin{table}[H]
\centering
\begin{tabular}{lcc}
\toprule
 & \textbf{VNS déterministe} & \textbf{VNS stochastique} \\
\midrule
Résultat ($\sum T_j$) & 10 & 10 \\
Approche & Exhaustive & Aléatoire + recherche locale \\
Reproductible & Oui & Oui (seed fixe) \\
Avantage & Garantit le meilleur voisin local & Explore plus largement \\
Adapté à & Petites instances & Grandes instances \\
\bottomrule
\end{tabular}
\caption{Comparaison des variantes VNS}
\end{table}

Sur cette instance (3 jobs, 3 machines), les deux variantes donnent le même résultat. Le VNS améliore la solution LTR de $\sum T_j = 22$ à $\sum T_j = 10$, mais n'atteint pas l'optimum ($\sum T_j = 3$, trouvé par le Branch and Bound). Cela illustre la nature approchée des métaheuristiques : le VNS reste piégé dans un optimum local que seule une méthode exacte peut dépasser.
